\subparagraph{Esempio attacco SPN}
Le Crittoanalisi lineari riguardano proprio il trovare correlazioni lineari delle S-box nel cifrario, combinandole per ottenere un'approssimazione di tutto il cifrario. 

\begin{table}[h]
\centering
\begin{tabular}{|c|*{16}{c}|}
\hline
 & \multicolumn{16}{c|}{$b$} \\
$a$ 
& 0 & 1 & 2 & 3 & 4 & 5 & 6 & 7 & 8 & 9 & A & B & C & D & E & F \\
\hline
0 & 16 & 8 & 8 & 8 & 8 & 8 & 8 & 8 & 8 & 8 & 8 & 8 & 8 & 8 & 8 & 8 \\
1 & 8 & 8 & 6 & 6 & 8 & 8 & 6 & 14 & 10 & 10 & 8 & 8 & 10 & 10 & 8 & 8 \\
2 & 8 & 8 & 6 & 6 & 8 & 8 & 6 & 6 & 8 & 8 & 10 & 10 & 8 & 8 & 2 & 10 \\
3 & 8 & 8 & 8 & 8 & 8 & 8 & 8 & 8 & 10 & 2 & 6 & 6 & 10 & 10 & 6 & 6 \\
4 & 8 & 10 & 8 & 6 & 6 & 4 & 6 & 8 & 8 & 6 & 8 & 10 & 10 & 4 & 10 & 8 \\
5 & 8 & 6 & 6 & 8 & 6 & 8 & 12 & 10 & 6 & 8 & 4 & 10 & 8 & 6 & 6 & 8 \\
6 & 8 & 10 & 6 & 12 & 10 & 8 & 8 & 10 & 8 & 6 & 10 & 12 & 6 & 8 & 8 & 6 \\
7 & 8 & 6 & 8 & 10 & 10 & 4 & 10 & 8 & 6 & 8 & 10 & 8 & 12 & 10 & 8 & 10 \\
8 & 8 & 8 & 8 & 8 & 8 & 8 & 8 & 8 & 6 & 10 & 10 & 6 & 10 & 6 & 6 & 2 \\
9 & 8 & 8 & 6 & 6 & 8 & 8 & 6 & 6 & 4 & 8 & 6 & 10 & 8 & 12 & 10 & 6 \\
A & 8 & 12 & 6 & 10 & 4 & 8 & 10 & 6 & 10 & 10 & 8 & 8 & 10 & 10 & 8 & 8 \\
B & 8 & 12 & 8 & 4 & 12 & 8 & 12 & 8 & 8 & 8 & 8 & 8 & 8 & 8 & 8 & 8 \\
C & 8 & 6 & 12 & 6 & 6 & 8 & 10 & 8 & 10 & 8 & 10 & 12 & 8 & 10 & 8 & 6 \\
D & 8 & 10 & 10 & 8 & 6 & 12 & 8 & 10 & 4 & 6 & 10 & 8 & 10 & 8 & 8 & 10 \\
E & 8 & 10 & 10 & 8 & 6 & 4 & 8 & 10 & 6 & 8 & 8 & 6 & 4 & 10 & 6 & 8 \\
F & 8 & 6 & 4 & 6 & 6 & 8 & 10 & 8 & 8 & 6 & 12 & 6 & 6 & 8 & 10 & 8 \\
\hline
\end{tabular}
\caption{Tabella delle approssimazioni lineari}
\label{tab:NL}
\end{table}


Si elencano le approssimazioni usate per le S-box, il pedice indica il round della S-box e l'indice della S-box presa in considerazione (da 1 a 4),
\begin{itemize}
    \item Per $S_{1,2}$ sia  $T_1 = U_{1,5} \oplus U_{1,7} \oplus U_{1,8} \oplus V_{1,6}$ con bias $1/4$;
    \item Per $S_{2,2}$ sia  $T_2 = U_{2,6} \oplus V_{2,6} \oplus V_{2,8}$ con bias $-1/4$;
    \item Per $S_{3,2}$ sia  $T_3 = U_{3,6} \oplus V_{3,6} \oplus V_{3,8}$ con bias $-1/4$;
    \item Per $S_{3,4}$ sia  $T_4 = U_{3,14} \oplus V_{3,14} \oplus V_{3,16}$ con bias $-1/4$.
\end{itemize}
Mentre $U,V$ sono le variabili che corrispondono all'input e output delle S-box. Assumendo che variabili $T_1, T_2, T_3, T_4$ siano indipendenti si può fare uso del Piling-Up Lemma, quindi $T_1\oplus T_2\oplus T_3\oplus T_4$ ha bias $\varepsilon = -1/32$. Inoltre le variabili sono state scelte in modo tale che l'output di $T_r$ sia l'input di $T_{r+1}$, perciò è possibile esprimere i bit risultanti in funzione del \emph{plaintext}

\begin{align*}
    T_1 &= U_{1,5} \oplus U_{1,7} \oplus U_{1,8} \oplus V_{1,6} = X_5 \oplus K_{1,5} \oplus X_7 \oplus X_8 \oplus K_{1,8} \oplus V_{1,6} \\
    T_2 &= U_{2,6} \oplus V_{2,6} \oplus V_{2,8} = V_{1,6} \oplus K_{2,6} \oplus V_{2,6} \oplus V_{2,8}\\
    T_3 &= U_{3,6} \oplus V_{3,6} \oplus V_{3,8} = V_{2,6} \oplus K_{3,6} \oplus V_{3,6} \oplus V_{3,8}\\
    T_4 &= U_{3,14} \oplus V_{3,14} \oplus V_{3,16} = V_{2,8} \oplus K_{3,14} \oplus V_{3,14} \oplus V_{3,16}
\end{align*}
si ottiene infine
\[
    T_1\oplus T_2\oplus T_3\oplus T_4 = X_5 \oplus X_7 \oplus X_8 \oplus V_{3,6} \oplus V_{3,8} \oplus V_{3,14} \oplus V_{3,16} \oplus K_{1,5} \oplus K_{1,7} \oplus K_{1,8} \oplus K_{2,6} \oplus K_{3,6} \oplus K_{3,14}. 
\]
sostituendo i valori di output con l'input del round precedente:
\begin{align*}
    V_{3,6} &= U_{4,6} \oplus K_{4,6}\\
    V_{3,8} &= U_{4,14} \oplus K_{4,14}\\
    V_{3,14} &= U_{4,8}\oplus K_{4,8}\\
    V_{3,16} &= U_{4,16} \oplus K_{4,16}
\end{align*}

ricaviamo:
\[
    X_5 \oplus X_7 \oplus X_8 \oplus U_{4,6} \oplus U_{4,8} \oplus U_{4,14} \oplus U_{4,16} \oplus K_{1,5} \oplus K_{1,7} \oplus K_{1,8} \oplus K_{2,6} \oplus K_{3,6} \oplus K_{3,14} \oplus K_{4,6} \oplus K_{4,8} \oplus K_{4,14} \oplus K_{4,16}.
\]

Separando i bit della chiave rimangono solamente variabili $U$ relative all'ultimo round 
\begin{equation}
    \label{chiave}
    K_{1,5} \oplus K_{1,7} \oplus K_{1,8} \oplus K_{2,6} \oplus K_{3,6} \oplus K_{3,14} \oplus K_{4,6} \oplus K_{4,8} \oplus K_{4,14} \oplus K_{4,16} 
\end{equation}

\begin{equation}
    \label{eq5}
    X_{5} \oplus X_{7} \oplus X_{8} \oplus U_{4,6} \oplus U_{4,8} \oplus U_{4,14} \oplus U_{4,16} 
\end{equation}


con cui si ricorda avere un bias $\varepsilon = -1/32$, di cui le uniche variabili sconosciute sono quelle della chiave \ref{chiave}. Se si cercassero di indovinare i bit della chiave con $T \approx 8000$ coppie di \emph{plaintext} e \emph{ciphertext}, si riuscirebbero a trovare i bit della chiave nel 5 round usati dalle S-box 2 e 4. Per una chiave candidata di 8 bit ci sono 256 possibili chiavi da verificare, per ogni coppia di \emph{plaintext} e \emph{ciphertext} e ogni chiave K viene verificato quando \ref{eq5} $= 0$, si incrementa un contatore associato alla chiave utilizzata, molti contatori saranno vicini al valore $T/2$, il contatore con valore più vicino al valore $T/32$ sarà associato alla chiave candidata corretta.
%Continua con la spiegazione, ma l'idea è questa

Generalmente per un attacco lineare c'è bisogno di $T \approx c \varepsilon^{-2}$ con $c$ costante e $\varepsilon$ bias.
